\documentclass[11pt]{article}

\usepackage[margin=1.1in]{geometry}
\usepackage{amsmath,amssymb,amsthm}
\usepackage{booktabs}
\usepackage{hyperref}

\newtheorem{theorem}{Theorem}[section]
\newtheorem{proposition}[theorem]{Proposition}
\newtheorem{lemma}[theorem]{Lemma}
\newtheorem{corollary}[theorem]{Corollary}
\theoremstyle{definition}
\newtheorem{definition}[theorem]{Definition}
\theoremstyle{remark}
\newtheorem{remark}[theorem]{Remark}

\newcommand{\A}{\mathbb A}
\newcommand{\F}{\mathbb F}
\newcommand{\Q}{\mathbb Q}
\newcommand{\Z}{\mathbb Z}
\newcommand{\PP}{\mathbb P}
\newcommand{\Jac}{\operatorname{Jac}}
\newcommand{\conv}{\operatorname{conv}}
\newcommand{\ord}{\operatorname{ord}}
\newcommand{\Proj}{\operatorname{Proj}}

\title{An Exact Computer-Assisted Exclusion at the Degree-\(125\)
Frontier of the Plane Jacobian Conjecture}
\author{Atharva Vaidya}
\date{25 July 2026}

\begin{document}
\maketitle

\begin{abstract}
Let \(k\) be an algebraically closed field of characteristic zero and let
\((P,Q)\in k[x,y]^2\) have nonzero constant Jacobian.  Combining the
public preprint reduction of Guccione--Guccione--Horruitiner--Valqui
with two exact eliminations, we prove that a noninvertible pair must satisfy
\[
   \max\{\deg P,\deg Q\}\geq 125.
\]
The first elimination is a uniform geometric obstruction for the
consecutive transverse-degree \((2,3)\) Newton family.  The second is an
exact computer-assisted elimination of the remaining case-c coefficient
system.  Its outer equation is exhausted by a five-point Hurwitz fiber;
the lower equations reduce to seven weighted modes; and a complete
finite-field weighted-projective fiber is certified empty.  Properness
then gives the characteristic-zero conclusion.  We specify the external
theorems used, the exact arithmetic trust base, and a cache-free
reproduction protocol.  This result is a bounded exclusion, not a proof
of the plane Jacobian conjecture.  An independent announcement of the
same numerical bound exists, and no priority claim is made.
\end{abstract}

\section{Statement and logical boundary}

\begin{definition}
A \emph{plane Keller pair} over a field \(k\) is a pair
\((P,Q)\in k[x,y]^2\) for which
\[
  \Jac(P,Q):=P_xQ_y-P_yQ_x\in k^\times.
\]
\end{definition}

\begin{theorem}[Bounded exclusion]\label{thm:main}
Let \(k\) be an algebraically closed field of characteristic zero.  If a
plane Keller pair \((P,Q)\) is not a polynomial automorphism, then
\[
  \max\{\deg P,\deg Q\}\geq 125.
\]
\end{theorem}

It is enough to prove the statement over \(\mathbb C\).  Indeed, the
coefficients of a proposed pair generate a finitely generated
characteristic-zero subfield \(k_0\), which embeds in \(\mathbb C\).
If the base-changed endomorphism of \(\mathbb C[x,y]\) were invertible,
then the original endomorphism of \(k_0[x,y]\) would be invertible by
faithfully flat descent.  Thus a noninvertible pair remains
noninvertible after a suitable embedding in \(\mathbb C\).

Theorem~\ref{thm:main} is not the two-dimensional Jacobian conjecture.
It excludes a finite degree range.  The passage from a general
counterexample of degree below \(125\) to the coefficient systems treated
here uses results from the public preprint of
Guccione--Guccione--Horruitiner--Valqui (GGHV)~\cite{GGHV2022}; those
results are external dependencies and are not reproved.
We use them with the algebraically closed characteristic-zero convention
inherited from the earlier papers cited by GGHV.

The precise imported chain is as follows.

\begin{enumerate}
\item GGHV Theorem 2.1, together with the exhaustive small-case table on
  its page 3, reduces a counterexample with maximum degree below \(125\),
  up to exchanging \(P,Q\), to the degree pair \((72,108)\) and one of
  the two corner cases \((8,28)\), \((9,27)\).
\item GGHV Proposition 4.1 and Corollary 5.7, the latter proved through
  Theorem 5.1, eliminate the \((9,27)\) corner configuration.
\item GGHV Proposition 4.3 reduces the remaining \((8,28)\) corner
  configuration to exactly the two Newton-polygon alternatives displayed
  in Section~\ref{sec:import}.
\end{enumerate}

Our contribution to this chain is to exclude both alternatives.  We do
not claim an independent derivation of the GGHV reduction, nor do we
claim priority for the resulting bound; see Remark~\ref{rem:priority}.

\section{Exact import of the two GGHV systems}\label{sec:import}

GGHV Proposition 4.3 produces a Laurent-polynomial pair satisfying
\[
 [P,Q]_{x,y}=x^2. \tag{2.1}
\]
Its case-c alternative has polygons
\[
\begin{aligned}
\Delta_P&=\conv\{(0,0),(1,0),(8,14),(8,16),(0,8)\},\\
\Delta_Q&=\conv\{(0,0),(2,1),(12,21),(12,24),(0,12)\}.
\end{aligned} \tag{2.2}
\]
The a/b alternative has the same polygons with \((0,8)\) and \((0,12)\)
omitted.  Equality of Newton polygons means that all displayed vertices
have nonzero coefficients.

Put
\[
   z=xy,\qquad h=xy^2.
\]
Then
\[
   x^ay^b=z^{\,2a-b}h^{\,b-a},\qquad
   \det\begin{pmatrix}2&-1\\-1&1\end{pmatrix}=1, \tag{2.3}
\]
so this is a unimodular change of exponent lattice.  Moreover,
\[
 \det\frac{\partial(z,h)}{\partial(x,y)}=h,\qquad
 x^2=\frac{z^4}{h^2},
\]
and, after an inessential sign convention,
\[
   \{P,Q\}_{z,h}=\frac{z^4}{h^3}. \tag{2.4}
\]

For a/b, enumerating every lattice point gives the five consecutive
blocks
\[
\begin{aligned}
P&=A(h)+zp_1(h)+z^2\frac{U(h)}h,\\
Q&=B(h)+zq_1(h)+z^2q_2(h)+z^3\frac{V(h)}h,             \tag{2.5}
\end{aligned}
\]
where
\[
\deg A=8,\quad \deg B=12,\quad \deg U=7,\quad \deg V=10.
\]
The fixed-pole coefficients can be normalized to
\[
 U(0)=V(0)=1,
\]
and the outer bracket row is
\[
 UV+2hUV'-3hU'V=1.                                   \tag{2.6}
\]

For case c the two extra vertices transform as
\[
 (0,8)\longmapsto(-8,8),\qquad
 (0,12)\longmapsto(-12,12).                           \tag{2.7}
\]
Thus case c has \(P\)-blocks \(z^{-8},\ldots,z^2\) and \(Q\)-blocks
\(z^{-12},\ldots,z^3\).  It is not a five-block system and must be
eliminated separately.

\begin{lemma}[Support inventory]\label{lem:support}
The polygons in (2.2) contain respectively \(61\) and \(125\) lattice
points.  After removing the two bracket-invisible additive constants and
the \(8+11\) outer coefficients of \(U,V\), the case-c lower system has
\[
  61+125-2-8-11=165
\]
coefficient positions.  The radial inventory used in the certificate is
exactly their image under (2.3), with no omitted or duplicated monomial.
\end{lemma}

\begin{proof}
The assertion is a finite convex-lattice enumeration.  The verifier
\texttt{current\_context/verify\_case\_c\_full\_certificate\_bridge.py}
reconstructs both polygons, enumerates their lattice points, applies
(2.3), and compares the result with the recurrence inventory.
\end{proof}

\section{The consecutive five-block alternative}\label{sec:fiveblock}

The argument used here applies uniformly, although only the scale
\(R=4\) instance is needed for Theorem~\ref{thm:main}.

\begin{theorem}[Consecutive \((2,3)\) obstruction]\label{thm:allscale}
For \(R\geq1\), suppose (2.5) satisfies (2.4), with
\[
\begin{gathered}
\deg A=2R,\qquad \deg B=3R,\\
\deg U=2R-1,\qquad\deg V=3R-2,\qquad U(0)=V(0)=1,
\end{gathered}
\]
and with the full vertices that make these degrees exact.  Then no such
five-block completion exists.
\end{theorem}

\begin{lemma}[Deficit-one form and coprimality]\label{lem:deficit-one}
In a completion from Theorem~\ref{thm:allscale}, \(p_1\ne0\), and there
are constants \(a,b\), not both zero, such that
\[
   p_1=(ah+b)U'-\frac a2U.                              \tag{3.1}
\]
Moreover,
\[
   \gcd(U,p_1)=1.                                      \tag{3.2}
\]
\end{lemma}

\begin{proof}
For a deficit-one kernel pair \((p_1,q_2)\), the combination
\(C=2Uq_2-3Vp_1\) recovers the pair: substituting the outer equation
into the linear bracket row gives
\[
\begin{aligned}
p_1&=hCU'-\frac12CU-\frac12hC'U,\\
q_2&=hCV'-\frac14CV-\frac34hC'V.
\end{aligned}
\]
The endpoint degree and divisibility conditions, compared coefficient by
coefficient from the two ends, leave exactly
\[
 C=1,\qquad
 C_*=\frac{U^2-U(0)^2-2U(0)U'(0)h}{h},
\]
as a basis.  Direct substitution gives
\[
 p_{1,1}=hU'-\frac12U,\qquad
 p_{1,C_*}
 =-2U(0)U'(0)p_{1,1}-U(0)^2U'.
\]
Their span is exactly (3.1).  This is a coefficientwise identity; the
script
\texttt{route\_bd\_universal\_radial\_kernel\_theorem.py} checks the
kernel exhaustion and both substitutions.  If \(p_1=0\), the
\(z^0\)-row of the bracket first gives \(q_1=0\), and its \(z^2\)-row
then gives \(A'q_3=0\), impossible because \(A'\ne0\) and \(q_3=V/h\ne0\).

The outer equation shows that every root of \(U\) is nonzero and simple
and that \(V\) does not vanish there.  Suppose that
\(U(\alpha)=p_1(\alpha)=0\).  Evaluating the bracket at \(h=\alpha\)
gives
\[
 -\bigl(A'+p_1'z+p_2'z^2\bigr)
  \bigl(q_1+2q_2z+3q_3z^2\bigr)
 =\alpha^{-3}z^4.
\]
Here
\[
 p_2'(\alpha)=\frac{U'(\alpha)}{\alpha}\ne0,\qquad
 q_3(\alpha)=\frac{V(\alpha)}{\alpha}\ne0.
\]
Thus both quadratic factors are monomials in \(z\), and
\[
 A'(\alpha)=p_1'(\alpha)=q_1(\alpha)=q_2(\alpha)=0.
\]
If \(a=0\), then \(p_1=bU'\) is nonzero at \(\alpha\).  If \(a\ne0\),
(3.1) first gives \(a\alpha+b=0\), and differentiation gives
\(p_1'(\alpha)=aU'(\alpha)/2\ne0\).  Both cases are contradictions.
\end{proof}

\begin{lemma}[Fixed-pole nonimmersion]\label{lem:nonimmersion}
Factor the boundary parametrization polynomially as
\[
 h\longmapsto \xi=H(h)\longmapsto(A_0(\xi),B_0(\xi)),
\]
where the second map is the normalization of its image.  At
\(\xi_0=H(0)\),
\[
 A_0'(\xi_0)=B_0'(\xi_0)=0.                            \tag{3.3}
\]
\end{lemma}

\begin{proof}
Suppose first that \(A_0'(\xi_0)\ne0\), and write the normalized image
locally as \(B_0=g(A_0)\).  Over \(\mathbb C(h)\), put \(N=Q-g(P)\).
The identity
\[
 \{P,N\}_{z,h}=z^4/h^3
\]
and the nonvanishing of \((A_0\circ H)'\) in \(\mathbb C(h)\) imply
\(N=O(z^5)\).  Comparing its first four Taylor coefficients gives
\[
\begin{aligned}
q_1&=g'p_1,\\
q_2&=g'p_2+\frac12g''p_1^2,\\
q_3&=g''p_1p_2+\frac16g'''p_1^3,\\
0&=\frac12g''p_2^2+\frac12g'''p_1^2p_2
  +\frac1{24}g''''p_1^4.                               \tag{3.4}
\end{aligned}
\]
Because \(p_2=1/h+O(1)\), \(q_3=1/h+O(1)\), and \(p_1\) is
polynomial, the third row forces \(g''p_1\) to be a unit at \(h=0\).
The fourth row then has an uncancelled pole of order two.

If \(A_0'(\xi_0)=0\) but \(B_0'(\xi_0)\ne0\), write
\(A_0=f(B_0)\).  The \(z^2\)-coefficient of \(P-f(Q)=O(z^5)\) is
\[
 p_2=f'q_2+\frac12f''q_1^2.
\]
Its right side is regular at \(h=0\), whereas \(p_2\) has a simple
pole.  This proves (3.3).
\end{proof}

\begin{lemma}[Root consumption]\label{lem:consumption}
If \(\beta\ne0\) is a root of \(p_1\) of multiplicity \(m\), then
\[
   \ord_\beta A'\geq2m-1.                              \tag{3.5}
\]
\end{lemma}

\begin{proof}
Lemma~\ref{lem:deficit-one} gives \(p_2(\beta)\ne0\).  First
\(A'(\beta)=0\).  Otherwise \(B=g(A)\) is a local graph and (3.4)
applies.  At \(p_1(\beta)=0\), its third row gives \(q_3(\beta)=0\).
The outer equation makes this zero simple.  The fourth row gives
\(g''(\beta)=0\), but differentiating the third row then gives
\(q_3'(\beta)=0\), a contradiction.

Complete the square locally:
\[
 \zeta=\sqrt{p_2}\left(z+\frac{p_1}{2p_2}\right),\qquad
 P=\zeta^2+H_1(h).
\]
The boundary \(z=0\) is \(\zeta=s(h)\), where
\(\ord_\beta s=m\), and
\[
 \{P,Q\}_{\zeta,h}=u(h)(\zeta-s(h))^4,\qquad u(\beta)\ne0.
\]
Writing \(Q=\sum_{j=0}^3\widetilde q_j(h)\zeta^j\), the coefficients
of \(\zeta^0,\zeta^2,\zeta^4\) are
\[
\begin{aligned}
H_1'\widetilde q_1&=-us^4,\\
2\widetilde q_1'-3H_1'\widetilde q_3&=6us^2,\\
2\widetilde q_3'&=u.                                   \tag{3.6}
\end{aligned}
\]
Put \(c=\ord_\beta H_1'\).  The last row says that
\(\widetilde q_3\) is a unit or has a simple zero.  If
\(c\le2m-2\), the first row gives
\[
\ord_\beta\widetilde q_1=4m-c\ge2m+2.
\]
In the middle row, \(\widetilde q_1'\) then has order at least
\(2m+1\), while \(H_1'\widetilde q_3\) has order at most \(2m-1\);
this cannot equal a term of order \(2m\).  Hence \(c\ge2m-1\).
Finally \(A'=H_1'+2ss'\), proving (3.5).
\end{proof}

\begin{lemma}[Marked chart exhaustion]\label{lem:charts}
The fixed-pole conditions and Lemma~\ref{lem:nonimmersion} leave only the
following two horizontal charts in the pulled-back parameter \(h\):
\[
\begin{array}{ll}
\text{\rm(i)}&
\ord_0(A-A(0))=2t+1,\quad
\ord_0(B-B(0))=3t+2,\quad
\ord_0p_1=t;\\[2mm]
\text{\rm(ii)}&
p_1(0)\ne0,\quad \ord_0 A'\ge2,
\end{array}                                             \tag{3.7}
\]
where in (ii) the normalized graph begins with a nonzero quadratic
term.  The opposite projection and every other analytic or fractional
chart are impossible.
\end{lemma}

\begin{proof}
We give the valuation exhaustion because the fixed poles make it much
shorter than a general Puiseux classification.  Write
\[
 m=\ord_0(A-A(0)),\qquad n=\ord_0(B-B(0)),\qquad
 r=\ord_0p_1.
\]
A nonzero tangent is impossible.  If no later fractional term cancels
its simple pole in \(g'p_2\), this follows already from the second row
of (3.4).  If a first fractional term of pulled-back order \(N>m\)
does cancel it, then necessarily
\(N-2m+2r=-1\).  Its three contributions to the fourth row have orders
\[
N-2m-2,\qquad -m-2,\qquad -N-2,
\]
so the last is uniquely lowest.  Suppose first that the horizontal graph
therefore begins with a nonintegral term
\(g(X)=cX^{n/m}+\cdots\).
The three valuations in the fourth row of (3.4) are
\[
 n-2m-2,\qquad n-3m+2r-1,\qquad n-4m+4r.               \tag{3.8}
\]
They form an arithmetic progression.  Unless
\(2r=m-1\), (3.8) has a unique lowest term.  With
\(2r=m-1\), the two terms in the third row of (3.4) have the same
valuation.  Their leading coefficients cannot cancel: writing \(T\) for
the normalized leading value of \(p_1^2h/(A-A(0))\), simultaneous
vanishing of the fourth-row leading coefficient and of
\(1+(n/m-2)T/6\), the third-row leading factor, would force
\(n/m=1/2\).  Since \(n/m>1\), the exact simple pole of \(q_3\) forces
\[
 2n=3m+1.
\]
Thus \(m=2t+1,n=3t+2,r=t\), which is (i).

The opposite projection is impossible.  Indeed, if
\(A=f(B)\), put \(s=\ord_0q_1\).  The fixed pole in
\[
 p_2=f'q_2+\frac12f''q_1^2
\]
forces \(m-2n+2s=-1\).  In the next Taylor row
\[
 0=f'q_3+f''q_1q_2+\frac16f'''q_1^3,
\]
the first two terms exceed the order of the last by respectively
\[
 \frac{3m-2n+1}{2}>0,\qquad m-n+\ord_0q_2+1>0.
\]
The last term is therefore uniquely lowest.

It remains to treat an analytic initial part.  A first analytic term of
degree at least three cannot supply the pole of \(q_3\).  Suppose instead
\[
 g(X)=c_2X^2+c_3X^3+cX^{N/m}+\cdots,\qquad c_2c\ne0.
\]
If \(2m<N<3m\), the fourth row of (3.4) is impossible as follows.
For \(r=0\), the third row forces \(N=3m-1\), after which the
\(g''''p_1^4\) term has unique order \(-m-1<-2\).  For \(r>0\), the
same row forces \(N=3m-3r-1\), unless \(m=2r+1\).  In the former case
the \(g''''p_1^4\) term is again uniquely lowest; in the exceptional
case all fractional terms have order above \(-2\), leaving
\(\frac12g''p_2^2\) uniquely at order \(-2\).

Consequently the first nonanalytic exponent lies above three.  The pole
of \(q_3\) then forces \(r=0\), and the only possible cancellation of
\(\frac12g''p_2^2\) in the fourth row is
\[
 N=4m-2.
\]
Hence \(m\ge3\) and \(\ord_0A'=m-1\ge2\), giving (ii).  If there is no
nonanalytic term, the image is a smooth graph; its normalization is an
isomorphism there, contrary to (3.3).

Finally, polynomial L\"uroth factorization gives
\(\mathbb C(A,B)=\mathbb C(H)\) for a polynomial \(H\).  If
\(e=\ord_0(H-H(0))\), all normalization orders above are multiplied by
\(e\).  In (i), \(2n=3m+1\) forces \(e=1\).  In (ii),
\(N=4m-2\) permits only \(e\mid2\), and the pulled-back inequalities
remain unchanged.  Thus no primitivity of \(h\mapsto(A(h),B(h))\) was
used, and the listed charts are exhaustive.
\end{proof}

\begin{proof}[Proof of Theorem~\ref{thm:allscale}]
In chart (i), the marked point contributes \(2t\) zeros to \(A'\).
If \(a\ne0\) in (3.1), then \(\deg p_1=2R-1\).  Applying
Lemma~\ref{lem:consumption} to every nonmarked root of \(p_1\) gives
\[
 \deg A'\ge2t+(2R-1-t)>2R-1,
\]
a contradiction.  If \(a=0\), then \(p_1=bU'\) has degree \(2R-2\)
and
\[
 \deg A'\ge2t+(2R-2-t).
\]
This is too large for \(t\ge2\).  For \(t=1\), equality forces
\[
 p_1=bU',\qquad A'=\lambda hU'
\]
with \(b,\lambda\ne0\).  The first lower bracket row gives
\[
 q_1=\frac b\lambda\frac{B'}h,\qquad \deg q_1=3R-2.
\]
Since \(U'(0)=0\), the other deficit-one partner is, up to a nonzero
scalar,
\[
 q_2=V'+\frac{V}{2h}-\frac{U}{2h},
\]
and has degree \(3R-3\).
The original \(z^1\)-row is
\[
h^2\left(
2\frac{U-1}{h}B'+p_1q_1'-p_1'q_1-2A'q_2
\right)+2hB'=0.                                       \tag{3.9}
\]
Inside the parentheses the four displayed degrees are
\[
5R-3,\qquad5R-5,\qquad5R-5,\qquad5R-4.
\]
Thus the first term in (3.9) is uniquely highest, a contradiction.

In chart (ii), every zero of \(p_1\) is nonmarked.  If \(a\ne0\),
Lemma~\ref{lem:consumption} gives
\[
\deg A'\ge2+(2R-1)>2R-1.
\]
If \(a=0\), it gives
\[
\deg A'\ge2+(2R-2)>2R-1.
\]
This excludes the second chart and completes the proof.
\end{proof}

\begin{corollary}\label{cor:ab}
The a/b alternative of GGHV Proposition 4.3 is empty.
\end{corollary}

\begin{proof}
Equations (2.3)--(2.6) identify the complete a/b support, including its
vertices and fixed poles, with Theorem~\ref{thm:allscale} at \(R=4\).
\end{proof}

\begin{remark}
The phrase ``all scale'' concerns the displayed consecutive five-block
family.  Known minimal-counterexample theory does not put every
possible degree or every corner chain into that family.  In particular,
Theorem~\ref{thm:allscale} by itself neither treats case c nor proves the
Jacobian conjecture.
\end{remark}

\section{The outer Hurwitz fiber in case c}\label{sec:hurwitz}

The outer blocks of case c still satisfy (2.6), with
\[
 \deg U=7,\qquad \deg V=10,\qquad U(0)=V(0)=1.          \tag{4.1}
\]
Associate the rational function
\[
   R(h)=h\frac{V(h)^2}{U(h)^3}.
\]
Equation (2.6) gives
\[
   R'(h)=\frac{V(h)}{U(h)^4}.                           \tag{4.2}
\]
The outer equation prevents a common root of \(U,V\).  At a root of
\(U\) it reads \(-3hU'V=1\), and at a root of \(V\) it reads
\(2hUV'=1\); hence all seven roots of \(U\) and all ten roots of \(V\)
are nonzero and simple.  Thus \(R\) has the first two ramification
partitions below.  If \(\lambda\) is the ratio of the leading
coefficients of \(hV^2\) and \(U^3\), then (4.2) has order \(-18\) at
infinity, so \(R-\lambda\) has order \(17\) in the local parameter
\(h^{-1}\).  Every other point above \(\lambda\) is unramified by
(4.2).  Degree \(21\) therefore leaves exactly four such points.
The resulting three-point passport is
\[
 (2^{10},1),\qquad(3^7),\qquad(17,1^4).                \tag{4.3}
\]

\begin{proposition}[Outer exhaustion]\label{prop:outer}
Up to the normalized source scaling, the geometric solution set of
(4.1)--(4.2) consists of five points.  It is represented by the five
embeddings of an irreducible quintic algebra
\[
   K=\Q[s]/(f(s)),\qquad s=u_7^{-1}.
\]
No additional solution lies on \(u_1=0\).
\end{proposition}

\begin{proof}
The exact symmetric-group character sum for (4.3) is \(5\).  Every
permutation triple counted is transitive: a component disjoint from the
17-cycle would lie among its four fixed sheets, where the involution and
the product of 3-cycles cannot restrict as mutual inverses.  Cover
automorphisms are trivial.  Indeed, they fix the unique simple zero and
the unique point of ramification index \(17\), so in the normalized
coordinate they have the form \(h\mapsto\lambda h\); the normalization
\(R(h)=h+O(h^2)\) gives \(\lambda=1\).

For an outer solution, normalize the third branch value by
\(T=R/\lambda\).  The scaling below sends
\(R(h)\) to \(R(\mu h)/\mu\), and hence sends \(T(h)\) to the
source-isomorphic cover \(T(\mu h)\).  Conversely, an isomorphism between
two normalized covers fixes the unique simple zero and the unique
\(17\)-fold point, so it is \(h\mapsto\mu h\); comparing
\(R_i'(0)=1\) shows that the two outer solutions differ by exactly this
action.  Thus outer scaling orbits inject into the Hurwitz fiber.
Exact rational reconstruction
produces an irreducible quintic \(f\) and five distinct solutions in the
chart \(u_1=1\), one at each geometric embedding.  That chart meets each
orbit with \(u_1\ne0\) once.  The five reconstructed orbits attain the
Hurwitz upper bound, so they exhaust the entire outer fiber and exclude
a separate orbit on \(u_1=0\).  The scripts
\texttt{route\_bd\_ab\_hurwitz\_count.py} and
\texttt{route\_bd\_ab\_outer\_lift.py} perform the character calculation,
irreducibility check, rational reconstruction, and coefficientwise
verification of (2.6).
\end{proof}

The GGHV case-c slice is commonly normalized by \(u_7=1\), whereas
the quintic above uses \(u_1=1\).  The bracket-preserving action
\[
 (P,Q)(z,h)\longmapsto
 \bigl(\lambda P(z,\lambda h),\lambda Q(z,\lambda h)\bigr)
\]
acts invertibly on every lower weighted coefficient.  On a quintic
representative, the equation \(u_7=1\) is
\(\lambda^7=u_7^{-1}=s\).  Since \(u_7\ne0\), it has seven distinct
solutions over the algebraic closure, and the action is free because
\(u_1\ne0\).  Thus the two normalizations differ only by an explicitly
audited invertible scaling; no outer or lower component is lost.

At the good prime
\[
   p=32003,                                             \tag{4.4}
\]
the primitive quintic has squarefree factorization
\[
 (s-26839)(s-16621)
 (s^3-11133s^2-11294s-6180).                           \tag{4.5}
\]
Thus the reduced outer fiber is the disjoint union of two rational points
and one degree-three point.  These factors account for all five geometric
outer points.

\section{Seven modes and the consistency scheme}\label{sec:modes}

Order the remaining blocks by radial deficit from outer degrees \(2\)
and \(3\).  At each stage the new coefficients occur linearly, while the
right-hand side depends polynomially on earlier blocks.

\begin{proposition}[Exact lower parametrization]\label{prop:modes}
On the complete outer algebra, the bounded homogeneous kernels at
successive deficits have dimensions
\[
   (2,2,2,1,0,0,\ldots,0).                              \tag{5.1}
\]
Consequently every solution of the \(165\)-coordinate case-c system
determines seven canonical parameters of weights
\[
   (1,1,2,2,3,3,4).                                    \tag{5.2}
\]
Every later coefficient is uniquely determined by these parameters, and
every retained cokernel row is a necessary equation for the original
system.
\end{proposition}

\begin{proof}
The universal kernel identity gives seven explicit modes and hence the
kernel lower bounds in (5.1).  Exact row reduction through all fifteen
deficits gives the matching upper bounds and vanishing of every later
kernel.  The calculation is repeated on the two linear factors and the
cubic factor in (4.5); nonzero maximal minors at good reduction certify
rank lower bounds, hence kernel upper bounds, while the explicit
characteristic-zero modes certify kernel lower bounds.  Since later
systems are affine-linear with source terms in earlier blocks, a failed
cokernel row cannot be repaired at a later deficit.
\end{proof}

After an invertible endpoint-coordinate change, the two first
nonautomatic rows, at deficit four, are nonzero scalar multiples of
\[
   \left(X_2-\frac{169}{48}X_0^2\right)^2.              \tag{5.3}
\]
For the transparent chart decomposition, set-theoretically we may
therefore impose
\[
   X_2=\frac{169}{48}X_0^2.                             \tag{5.4}
\]
The surviving compatibility rows then have the following inventory:
\[
\begin{array}{c|cccc}
\text{weight}&5&6&7&8\\ \midrule
\text{number of rows}&1&3&5&6.
\end{array}                                             \tag{5.5}
\]
All equations are homogeneous for the weights in (5.2).  Let
\(I_8^{\mathrm{orig}}\) denote the ideal generated by the original
consistency rows through deficit eight, retaining the squared rows (5.3).
For the characteristic-zero transport, let
\(I_8^{\mathrm{rad}}\) instead denote the homogeneous ideal generated by
\[
 X_2-\frac{169}{48}X_0^2
\]
and the exact rows in (5.5) after this substitution.  In characteristic
zero, and at the audited good factors above \(32003\), the two ideals
have the same field-valued projective zero set, although their scheme
structures need not agree.

\section{The complete special-fiber certificate}\label{sec:certificate}

\begin{theorem}[Finite-field certificate]\label{thm:finite}
For each of the three factors in (4.5), exact standard-basis reduction in
the corresponding field gives
\[
   X_i^{32}\in I_8^{\mathrm{orig}}\qquad(0\le i\le6),  \tag{6.1}
\]
and
\[
   \dim_{\F_{p^e}}
   \F_{p^e}[X_0,\ldots,X_6]/I_8^{\mathrm{orig}}=380,
   \qquad e=1,1,3.                                     \tag{6.2}
\]
In particular,
\[
 \Proj\bigl(
 \F_{p^e}[X_0,\ldots,X_6]/I_8^{\mathrm{orig}}
 \bigr)=\varnothing
 \quad(e=1,1,3).                                       \tag{6.3}
\]
\end{theorem}

\begin{proof}
All arithmetic is exact.  Singular reduces each of the seven monomials
in (6.1) to zero against an exact finite-field standard basis.  Hence the
radical of \(I_8^{\mathrm{orig}}\) contains the weighted irrelevant ideal
\((X_0,\ldots,X_6)\), proving (6.3).  The independently computed quotient
dimension (6.2) is a consistency check on all three standard-basis runs.
\end{proof}

Separately, the radical-branch scheme defined by
\(I_8^{\mathrm{rad}}\) admits the exhaustive chart cover
\[
D(L),\quad V(L)\cap D(A),\quad
V(L,A)\cap D(B),\quad V(L,A,B).                         \tag{6.4}
\]
The weight-five row is linear in \(X_6\) with coefficient \(L\); after
\(L=0\), its coefficients of \(X_5,X_4\) are \(A,B\).  The first two
charts give unit ideals at good reduction.  The third is empty by a
nonzero resultant of two quadratics.  On the fourth, a nonzero
\(3\times3\) determinant forces \(X_4=X_5=0\), deficit seven is automatic,
and the deficit-eight rows are nonzero multiples of \(X_6^2\).
The determinant and resultant statements are also checked directly in
the exact quintic field \(K\).  It is essential to use all four charts:
emptiness of \(D(L)\) alone would not prevent a generic point from
specializing to \(L=0\).  This four-chart calculation, rather than a
coefficientwise lift of \(I_8^{\mathrm{orig}}\), is the special fiber
used below.

\begin{theorem}[Characteristic-zero lift]\label{thm:lift}
The weighted-projective radical-branch consistency scheme has empty
characteristic-zero geometric fiber.
\end{theorem}

\begin{proof}
The outer equation, the endpoint change, the square root equation (5.4),
and every substituted row in (5.5) are constructed over \(K\).  After
inverting the finitely many denominators, their coefficients are integral
at the primes above \(32003\).  An independent transport check reduces
every one of these exact radical-branch rows,
coefficient-for-coefficient and in its original order, to the rows used
in the four-chart calculation (6.4).

These weighted homogeneous equations define a closed subscheme
\[
 \mathcal X^{\mathrm{rad}}\subset\PP_R(1,1,2,2,3,3,4)
\]
over the corresponding localized integral coefficient ring \(R\).
Weighted projective space is proper over \(R\), for example by passing to
a Veronese subring, and therefore
\(\mathcal X^{\mathrm{rad}}\to\operatorname{Spec}R\) is proper.  If its
characteristic-zero geometric generic fiber were nonempty, the proper
image would be a closed subset containing the generic point of the
integral base, and hence would contain every closed point.  This
contradicts the complete empty radical-branch fiber in (6.4).  No
flatness assumption is used.
\end{proof}

\begin{remark}[Secondary characteristic-zero certificates]
The two special charts in (6.4) have direct determinant/resultant
certificates over \(K\).  On \(L=0,A\ne0\), ordinary exact standard-basis
reduction gives
\[
 A^8\in
 (L,F_5,F_{6,0},F_{6,1},F_{6,2},F_{7,0}).
\]
There is a logically exact proposed direct route on \(L\ne0\): homogenize
the affine chart ideal and use Singular \texttt{modStd(...,1)} to test a
power of the homogenizing variable.  That run was stopped after a bounded
audit window and no successful power-reduction output is archived.
Accordingly it is not a certificate and is not used in
Theorem~\ref{thm:lift}.
\end{remark}

\section{Excluding the cone origin}\label{sec:origin}

The finite-field certificate leaves only the affine cone origin.  That
origin is not a genuine case-c point.

\begin{proposition}\label{prop:origin}
Every case-c solution with the Newton polygons (2.2) determines a
nonzero point of
\(\PP(1,1,2,2,3,3,4)\).
\end{proposition}

\begin{proof}
The required opposite vertices occur at
\[
 [z^{-8}h^8]P\ne0,\qquad[z^{-12}h^{12}]Q\ne0.           \tag{7.1}
\]
Their radial deficits are respectively \(10\) and \(15\).  If all seven
modes vanish, the first four lower blocks vanish.  Every later
inhomogeneous source is then zero; by (5.1), the corresponding linear
operator has zero kernel.  Induction through deficits \(5,\ldots,15\)
forces all later lower coefficients to vanish, including both
coefficients in (7.1).  This contradicts equality with the required
Newton polygons.  The two omitted additive constants are
bracket-invisible and do not contribute to (7.1).
\end{proof}

\begin{corollary}[Case-c elimination]\label{cor:casec}
The coefficient system in GGHV Proposition 4.3(1) has no solution over
an algebraically closed field of characteristic zero.
\end{corollary}

\begin{proof}
By Propositions~\ref{prop:outer} and~\ref{prop:modes}, any solution gives
a nonzero point satisfying the original weighted-projective consistency
rows.  The exact square identity (5.3) then puts that point on the
radical-branch scheme.  Theorem~\ref{thm:lift} says that no such point
exists.
\end{proof}

\begin{proof}[Proof of Theorem~\ref{thm:main}]
The external GGHV reduction listed in Section~1 leaves, below degree
\(125\), only the two alternatives in Section~\ref{sec:import}.
Corollary~\ref{cor:ab} excludes a/b, and
Corollary~\ref{cor:casec} excludes case c.
\end{proof}

\section{Trust base and reproduction}\label{sec:reproduce}

The proof separates mathematical reductions from machine certificates.
Its trusted computational base is:

\begin{itemize}
\item exact integer, rational, algebraic-number, and finite-field
  arithmetic;
\item SymPy \(1.14.0\) exact polynomial arithmetic~\cite{SymPy};
\item Singular \(4.4\) or newer and its exact finite-field
  standard-basis algorithm~\cite{Singular}; and
\item for interpreting attempted direct homogeneous
  characteristic-zero runs only, the documented exactness mode of
  \texttt{modstd.lib}.  No such run is a premise of the proof.
\end{itemize}

Floating-point root selection is never used.  The nonhomogeneous
\texttt{modStd} unit outputs, which Singular documents only as
high-probability in the relevant global setting, are not premises of
Theorem~\ref{thm:main}.

From the repository root~\cite{Archive}, create the pinned Python
environment and run the quick interface audit:
\begin{verbatim}
python3 -m venv .venv
.venv/bin/pip install -r requirements.txt
.venv/bin/python \
  current_context/verify_case_c_full_certificate_bridge.py
\end{verbatim}
The complete finite-field weighted-projective certificate is:
\begin{verbatim}
.venv/bin/python \
  current_context/verify_case_c_full_certificate_bridge.py \
  --run-special-fiber
\end{verbatim}
The modular square branch and its two generic charts are:
\begin{verbatim}
.venv/bin/python route_bd_ab_hurwitz_count.py
.venv/bin/python route_bd_case_c_n3_hurwitz_bridge.py
.venv/bin/python route_bd_case_c_n3_generic_charts.py
\end{verbatim}
The exact characteristic-zero square, row transport, and special-chart
scalar certificates are:
\begin{verbatim}
.venv/bin/python scratch_case_c_n3_special_charts_Q.py
\end{verbatim}
For completeness, the commands below launch an optional direct
homogeneous generic-chart experiment:
\begin{verbatim}
.venv/bin/python scratch_case_c_n3_generic_charts_Q.py \
  --skip-modular-crosscheck --chart L_NE_0 --one-certificate
.venv/bin/python scratch_case_c_n3_generic_charts_Q.py \
  --skip-modular-crosscheck --chart L_EQ_0_A_NE_0 --one-certificate
\end{verbatim}
The \(L\ne0\) run was not completed in the present audit, and no
\texttt{HPOWER} output from this route is claimed.  These commands are
not part of the proof reproduction protocol.
The all-scale calculation and its symbolic identities are checked by:
\begin{verbatim}
.venv/bin/python route_bd_universal_radial_kernel_theorem.py
.venv/bin/python route_bd_allscale_marked_cusp_obstruction.py
.venv/bin/python current_context/verify_gghv_allscale_scope.py
\end{verbatim}

Three files under \texttt{tmp/} are runtime caches:
\begin{verbatim}
tmp/case_c_n3_generic_charts_Q.pkl
tmp/case_c_n3_special_charts_Q.pkl
tmp/case_c_n3_special_recurrence_Q.pkl
\end{verbatim}
They are accelerators, not premises.  A publication replay must move
them out of \texttt{tmp/}, reconstruct from scratch, and archive:

\begin{enumerate}
\item standard output and standard error;
\item Python, SymPy, Singular, and \texttt{modstd.lib} versions;
\item SHA-256 hashes of every verifier and log;
\item the factorization (4.5), all exact-to-modular row comparisons, and
  all seven reductions (6.1) on every factor; and
\item regenerated cache hashes as diagnostics only.
\end{enumerate}
The canonical replay driver is
\begin{verbatim}
publication_artifacts/degree-125-bound/cleanroom-2026-07-25/replay.sh
\end{verbatim}
and must be invoked from the repository root after the three caches have
been moved aside.

\section{Limitations and review status}

The theorem does not address the admissible corner chains at degree
\(125\) and above.  It does not provide the global exhaustion theorem
that would be required to deduce the plane Jacobian conjecture from
Theorem~\ref{thm:allscale}.  It also does not reprove the GGHV
small-degree classification.

The mathematical certificate has three deliberately redundant audit
layers: the global nilpotence test (6.1), the four-chart decomposition
(6.4), and direct exact-\(K\) scalar certificates on the special charts.
The accompanying publication bundle includes a cache-free replay, its
archived log, and source hashes.  Before journal submission, the
recommended remaining check is an independent line-by-line human audit
of the imported GGHV implications.

\begin{remark}[Independent announcement]\label{rem:priority}
A MathOverflow answer posted in 2026 announces an independent
computer-assisted exclusion giving the same degree-\(125\) lower bound
and states that a write-up is in preparation~\cite{MO2026}.  The present
manuscript does not rely on that announcement and makes no claim of
priority or uniqueness.
\end{remark}

\paragraph{Disclosure.}
AI systems assisted with symbolic-search orchestration, code review, and
draft preparation.  The named author is responsible for the mathematical
claims, the final source, and independent verification of the
certificates.

\begin{thebibliography}{9}

\bibitem{GGHV2022}
J.~A. Guccione, J.~J. Guccione, R. Horruitiner, and C. Valqui,
\emph{Increasing the degree of a possible counterexample to the Jacobian
Conjecture from 100 to 108},
preprint, arXiv:2204.14178, 2022,
\url{https://arxiv.org/abs/2204.14178}.

\bibitem{SymPy}
A. Meurer et al.,
\emph{SymPy: symbolic computing in Python},
PeerJ Computer Science \textbf{3} (2017), e103,
\url{https://doi.org/10.7717/peerj-cs.103}.

\bibitem{Singular}
W. Decker, G.-M. Greuel, G. Pfister, and H. Sch\"onemann,
\emph{Singular 4-4: A computer algebra system for polynomial
computations},
\url{https://www.singular.uni-kl.de}.

\bibitem{Archive}
A. Vaidya,
\emph{Mathematics research archive: exact Jacobian-conjecture
certificates},
\url{https://github.com/AtharvaVaidya/mathematics-research}.

\bibitem{MO2026}
J.~A. Guccione,
\emph{Answer to ``The simplest case of Jacobian conjecture''},
MathOverflow, 2026,
\url{https://mathoverflow.net/questions/513413/the-simplest-case-of-jacobian-conjecture/513458}.

\end{thebibliography}

\end{document}
